MPRA等功能分析技术与不断发展的序列-活性模型，已推动顺式调控元件（CRE）研究从事后注释走向预测性序列设计~\citep{Melnikov2012,Arnold2013,Tewhey2016,Inoue2017,Gordon2020,Kircher2019,Avsec2021,Taskiran2024,Gosai2024}。下一步的工程问题并不只是找到更高分的序列，而是把模型预测转化为少量可追溯、可合成、可比较并能够在实验中迭代的构建体。单独的模型得分无法说明一次编辑能否跨预测器保持，也无法说明其是否便于实际构建。CRE特别适合连接生物学与工程：其转录因子结合位点及周围序列上下文编码基因表达的时空特异性~\citep{Wasserman2004,ENCODE2012,Thurman2012,Roadmap2015,Sahu2022}，而其长度适合直接合成和报告基因测量。因此，可靠的设计流程既可用于检验细胞类型选择性调控语法，也可服务于合成生物学、基因递送和疾病相关非编码变异的功能优先排序。

现有调控序列设计主要侧重从头生成或直接最大化单一预测器的得分~\citep{Vaishnav2022,Almeida2024,Gosai2024}。这类方法能够找到高分序列，却可能大幅改写构建体，甚至完全脱离天然CRE环境。由此产生的候选更难归因于具体核苷酸变化，增加合成和筛选负担，也难以为下一轮设计提供清晰依据。对于实验和工程流程，更有价值的输出是一个紧凑的可追溯面板：保留亲本骨架，只改变明确数量的核苷酸，并能逐条说明设计依据。

我们将这一问题定义为\emph{最小编辑}：在几乎完整保留亲本CRE的前提下提高细胞类型选择性。亲本序列保留天然上下文，固定编辑预算同时约束设计距离、构建负担以及需要解释的因果变化数量。因此，每条变异都是对既有调控骨架的可控扰动，而不是缺乏参照的合成序列。所得替换明确、可审计，并可直接转入寡核苷酸合成和报告基因筛选，使最小编辑成为模型排序与实验迭代之间的实际接口。

但编辑数量少并不等于设计可靠。单一预测器的搜索可能偏好利用模型特有行为的序列，即使该预测器在天然留出序列上表现良好~\citep{Whalen2022,deBoer2024,Ovadia2019,Rafi2025}。贪婪优化保留单一目标下即时增益最大的替换；更宽的束搜索虽能探索更多路径，但当同一模型同时负责生成和排序时，风险仍然存在。因此，尚未解决的核心问题是候选选择：如何筛出一组紧凑、可解释的编辑，使预测效应能够在独立训练、架构不同的模型之间保持，同时将不确定性和序列域行为控制在合理范围内。只有在这一步，预测模型才真正转化为工程决策系统，而不只是分数生成器。

为解决这一选择问题，我们开发了SafeEdit-CRE（图~\ref{fig:workflow}）。已发表的Malinois预测器~\citep{Gosai2024}提供主要的细胞类型特异性目标；在约束束搜索过程中，架构差异的残差-扩张CNN集成负责评审候选，经过验证集校准的不确定性以及链一致性、GC偏移、同聚物和6-mer天然性等防护指标用于过滤不稳定或序列特征异常的编辑。第二个独立训练的多核集成完全排除在候选生成之外，仅用于最终跨模型评估。搜索、评审和评估的程序化分离，使模型一致性成为明确的设计准则，而不是事后检查。1、5、10和20个核苷酸的固定预算，则使设计自由度能够被定量比较，而不会将最小编辑与无约束生成混为一谈。

我们将SafeEdit-CRE与随机替换、迭代贪婪编辑以及具有相同束宽和扩展计划的计算量匹配主模型束搜索进行比较。评估包括三种细胞中600个留出天然CRE亲本的跨模型泛化、90个非重叠亲本上的锁定九准则设计质量基准、按细胞类型和编辑预算的分层分析，以及完整重搜索消融。SafeEdit-CRE提高了相对于贪婪编辑的跨模型特异性裕度，在与计算量匹配的束搜索相比保持预测特异性高度吻合的同时降低了集成不确定性，并将锁定设计审计通过率从38.8\%提高至60.6\%。最终Tier A文库包含覆盖三种细胞类型和四种编辑预算的74条变异。因此，本研究不仅严格检验了跨模型最小编辑，还交付了一项可直接使用的工程成果：一个紧凑、可追溯、可进入报告基因筛选的候选面板，能够减少需要订购和筛选的构建体数量，同时保留足够的序列上下文来解释细胞类型选择性的变化。

\begin{figure}[t]
  \centering
  \includegraphics[width=\textwidth]{fig1_workflow.pdf}
  \caption{\textbf{SafeEdit-CRE研究设计。} 天然CRE亲本在固定替换预算下进行编辑。Malinois提供主要特异性目标，残差-扩张CNN集成在约束束搜索期间提供评审。经过验证集校准的不确定性和序列域行为控制防范不可信编辑。选定序列随后由第二个独立训练的多核集成评估，用于跨模型泛化评估。}
  \label{fig:workflow}
\end{figure}
